\section{等距浸入}\label{sec:isometric-immersions}

许多 Riemann 流形是其他 Riemann 流形的子流形, 并取诱导度量 (例如 $S^n\subset\mathbb{R}^{n+1}$). 本节分析这两个流形的曲率之间的关系.

设 $f:N\to M$ 是一个 $n$ 维流形 $N$ 到 $m$ 维流形 $M$ 的浸入. 由第~\ref{chap:differentiable-manifolds}~章第~\RGref{src:immersion-section-in-chapter-one}~节可知, 对每一点 $p\in N$, 都存在 $p$ 的邻域 $V\subset N$, 使 $f$ 在 $V$ 上是到其像的嵌入. 因而 $f(V)$ 是 $M$ 的子流形. 为简化记号, 将 $V$ 与 $f(V)$ 等同, 把 $f$ 当作包含映射继续讨论. 设 $\langle\cdot,\cdot\rangle$ 是 $M$ 上的 Riemann 度量, $\langle\!\langle\cdot,\cdot\rangle\!\rangle$ 是由 $f$ 在 $N$ 上诱导的度量 (因而称 $f$ 为\textbf{等距浸入}). 对每个 $p\in V$, 切空间 $T_pM$ 可以分解为
\[
  T_pM=T_pN\oplus(T_pN)^\perp.
\]
因此 $T_pM$ 的每个元素 $v$ 都可唯一写成 $v=v^\top+v^\perp$, 其中 $v^\top\in T_pN$ 是 $v$ 的\textbf{切向分量}, $v^\perp\in(T_pN)^\perp$ 是 $v$ 的\textbf{法向分量}. 设 $\widetilde\nabla$ 和 $\nabla$ 分别是 $(M,\langle\cdot,\cdot\rangle)$ 和 $(N,\langle\!\langle\cdot,\cdot\rangle\!\rangle)$ 的 Levi–Civita 联络. 设 $X,Y$ 是 $V\subset N$ 上的两个向量场, $\widetilde X,\widetilde Y$ 是 $X,Y$ 到 $V$ 的一个邻域 $W\subset M$ 的两个延拓. 利用 Koszul 公式, 容易验证
\[
  \nabla_XY=(\widetilde\nabla_{\widetilde X}\widetilde Y)^\top
\]
[参见第~\ref{chap:riemannian-manifolds}~章的 Exercise~\ref{ex:levi-civita-connection}(\ref{ex:induced-connection-and-projection})]. 定义 $N$ 的\textbf{第二基本形式}为
\[
  B(X,Y):=\widetilde\nabla_{\widetilde X}\widetilde Y-\nabla_XY.
\]
注意, 这个映射是良定义的, 即它不依赖于 $X,Y$ 的延拓 $\widetilde X,\widetilde Y$ [参见 Exercise~\ref{ex:isometric-immersions}(\ref{ex:second-fundamental-form-properties})]. 此外, 它双线性且对称, 对每个 $p\in V$, $B(X,Y)_p\in(T_pN)^\perp$ 只依赖于 $X_p$ 和 $Y_p$ 的值.

利用第二基本形式, 可以为每个向量 $n_p\in(T_pN)^\perp$ 定义一个对称双线性映射 $H_{n_p}:T_pN\times T_pN\to\mathbb{R}$:
\[
  H_{n_p}(X_p,Y_p)=\langle B(X_p,Y_p),n_p\rangle.
\]
相应的二次型通常称为 \textbf{$f$ 在 $p$ 处沿向量 $n_p$ 的第二基本形式}.

最后, 由于 $H_{n_p}$ 是双线性的, 存在线性映射 $S_{n_p}:T_pN\to T_pN$, 对所有 $X_p,Y_p\in T_pN$ 满足
\[
  \langle\!\langle S_{n_p}(X_p),Y_p\rangle\!\rangle
  =H_{n_p}(X_p,Y_p)=\langle B(X_p,Y_p),n_p\rangle.
\]
容易验证, 此线性映射由下式给出:
\[
  S_{n_p}(X_p)=-(\widetilde\nabla_{\widetilde X}n)_p^\top,
\]
其中 $n$ 是 $n_p$ 的一个垂直于 $N$ 的局部延拓. 事实上, 由于在 $N$ 上有 $\langle\widetilde Y,n\rangle=0$ 且 $\widetilde X$ 与 $N$ 相切, 所以在 $N$ 上有
\begin{align*}
  \langle\!\langle S_n(X),Y\rangle\!\rangle
  &=\langle B(X,Y),n\rangle
  =\langle\widetilde\nabla_{\widetilde X}\widetilde Y-\nabla_XY,n\rangle\\
  &=\langle\widetilde\nabla_{\widetilde X}\widetilde Y,n\rangle
  =\widetilde X\cdot\langle\widetilde Y,n\rangle
    -\langle\widetilde Y,\widetilde\nabla_{\widetilde X}n\rangle\\
  &=\langle-\widetilde\nabla_{\widetilde X}n,\widetilde Y\rangle
    =\langle\!\langle-(\widetilde\nabla_{\widetilde X}n)^\top,Y\rangle\!\rangle.
\end{align*}
因此, 对所有 $Y_p\in T_pN$,
\[
  \langle\!\langle S_{n_p}(X_p),Y_p\rangle\!\rangle
    =\langle\!\langle-(\widetilde\nabla_{\widetilde X}n)_p^\top,Y_p\rangle\!\rangle.
\]

\begin{RiemannExample}\label{ex:principal-curvatures-of-hypersurface}
设 $N$ 是 $M$ 中的一个\textbf{超曲面}, 即 $\dim N=n$ 且 $\dim M=n+1$. 取点 $p\in V$ ($V$ 是 $N$ 中 $f$ 为嵌入的一个邻域), 并取 $p$ 处垂直于 $N$ 的单位向量 $n_p$. 由于线性映射 $S_{n_p}:T_pN\to T_pN$ 对称, $T_pN$ 存在一组由特征向量 $\{(E_1)_p,\ldots,(E_n)_p\}$ 构成的标准正交基 (称为 $p$ 处的\textbf{主方向}), 对应于 $S_{n_p}$ 的实特征值 $\lambda_1,\ldots,\lambda_n$ (称为 $p$ 处的\textbf{主曲率}). 映射 $S_{n_p}$ 的行列式 (等于乘积 $\lambda_1\cdots\lambda_n$) 称为 $f$ 的 \textbf{Gauss 曲率}, 而 $H:=\frac1n\operatorname{tr}S_{n_p}=\frac1n(\lambda_1+\cdots+\lambda_n)$ 称为 $f$ 的\textbf{平均曲率}. 当 $n=2$ 且 $M=\mathbb{R}^3$ 带 Euclidean 度量时, $f$ 的 Gauss 曲率实际上就是第~\ref{sec:curvature}~节所定义的 $N$ 的 Gauss 曲率 (参见 Example~\ref{ex:theorema-egregium}).
\end{RiemannExample}

\begin{RiemannExample}\label{ex:derivative-of-gauss-map}
若上一例中 $M=\mathbb{R}^{n+1}$ 带 Euclidean 度量, 可以定义取值于单位球面的 \textbf{Gauss 映射} $g:V\subset N\to S^n$, 它给每个点 $p\in V$ 指派单位法向量 $n_p$. 由于 $n_p$ 垂直于 $T_pN$, 可以将切空间 $T_pN$ 与 $T_{g(p)}S^n$ 等同, 并得到一个良定义的映射 $(\mathrm{d}g)_p:T_pN\to T_pN$. 取满足 $c(0)=p$ 和 $\dot c(0)=X_p\in T_pN$ 的曲线 $c:I\to N$, 则
\[
  (\mathrm{d}g)_p(X_p)=\left.\frac{\mathrm{d}}{\mathrm{d}t}(g\circ c)\right|_{t=0}
  =\left.\frac{\mathrm{d}}{\mathrm{d}t}n_{c(t)}\right|_{t=0}
  =(\widetilde\nabla_{\dot c}n)_p,
\]
其中用到了 $\widetilde\nabla$ 是 Euclidean 度量的 Levi–Civita 联络这一事实. 然而, 由于 $\lVert n\rVert=1$, 有
\[
  0=\dot c(t)\cdot\langle n,n\rangle=2\langle\widetilde\nabla_{\dot c}n,n\rangle,
\]
从而
\[
  (\mathrm{d}g)_p(X_p)=(\widetilde\nabla_{\dot c}n)_p
    =(\widetilde\nabla_{\dot c}n)_p^\top=-S_{n_p}(X_p).
\]
于是 $p$ 处 Gauss 映射的导数为 $(\mathrm{d}g)_p=-S_{n_p}$.
\end{RiemannExample}

现在将 $N$ 与 $M$ 的曲率联系起来.

\begin{RiemannProposition}\label{prop:sectional-curvature-of-isometrically-immersed-submanifold}
设 $p$ 是 $N$ 中一点, $X_p,Y_p$ 是 $T_pN\subset T_pM$ 中两个线性无关的向量, $\Pi\subset T_pN\subset T_pM$ 是由这些向量生成的 2 维子空间. 设 $K^N(\Pi)$ 和 $K^M(\Pi)$ 分别表示 $N$ 和 $M$ 中相应的截面曲率. 则
\[
  K^N(\Pi)-K^M(\Pi)=
  \frac{\langle B(X_p,X_p),B(Y_p,Y_p)\rangle-\lVert B(X_p,Y_p)\rVert^2}
       {\lVert X_p\rVert^2\lVert Y_p\rVert^2-\langle X_p,Y_p\rangle^2}.
\]
\end{RiemannProposition}

\begin{RiemannProof}
注意, 右端只依赖于 $\Pi$, 所以不失一般性可假设 $\{X_p,Y_p\}$ 标准正交. 设 $X,Y$ 是 $X_p,Y_p$ 的局部延拓, 定义在 $p$ 在 $N$ 中的邻域上, 与 $N$ 相切且也标准正交. 设 $\widetilde X,\widetilde Y$ 是 $X,Y$ 到 $p$ 在 $M$ 中的邻域的延拓. 此外, 考虑也定义在 $p$ 在 $M$ 中的邻域上的标架场 $\{E_1,\ldots,E_{n+k}\}$, 使 $E_1,\ldots,E_n$ 与 $N$ 相切, 且在 $N$ 上有 $E_1=X$ 和 $E_2=Y$, 而 $E_{n+1},\ldots,E_{n+k}$ 垂直于 $N$ ($n+k=m$). 由于 $B(X,Y)$ 垂直于 $N$, 有
\[
  B(X,Y)=\sum_{i=1}^k\langle B(X,Y),E_{n+i}\rangle E_{n+i}
    =\sum_{i=1}^kH_{E_{n+i}}(X,Y)\,E_{n+i}.
\]
另一方面,
\begin{align*}
  K^N(\Pi)-K^M(\Pi)
   &=-R^N(X_p,Y_p,X_p,Y_p)+R^M(\widetilde X_p,\widetilde Y_p,\widetilde X_p,\widetilde Y_p)\\
   &=\bigl\langle\bigl(-\nabla_X\nabla_YX+\nabla_Y\nabla_XX+\nabla_{[X,Y]}X\\
   &\qquad{}+\widetilde\nabla_{\widetilde X}\widetilde\nabla_{\widetilde Y}\widetilde X
              -\widetilde\nabla_{\widetilde Y}\widetilde\nabla_{\widetilde X}\widetilde X
              -\widetilde\nabla_{[\widetilde X,\widetilde Y]}\widetilde X\bigr)_p,Y_p\bigr\rangle\\
   &=\bigl\langle\bigl(-\nabla_X\nabla_YX+\nabla_Y\nabla_XX
      +\widetilde\nabla_{\widetilde X}\widetilde\nabla_{\widetilde Y}\widetilde X
      -\widetilde\nabla_{\widetilde Y}\widetilde\nabla_{\widetilde X}\widetilde X\bigr)_p,Y_p\bigr\rangle,
\end{align*}
其中用到了 $\widetilde\nabla_{[\widetilde X,\widetilde Y]}\widetilde X-\nabla_{[X,Y]}X$ 垂直于 $N$ [参见 Exercise~\ref{ex:isometric-immersions}(\ref{ex:second-fundamental-form-properties})]. 然而, 由于在 $N$ 上
\begin{align*}
  \widetilde\nabla_{\widetilde Y}\widetilde\nabla_{\widetilde X}\widetilde X
  &=\widetilde\nabla_{\widetilde Y}(B(X,X)+\nabla_XX)\\
  &=\widetilde\nabla_{\widetilde Y}\left(\sum_{i=1}^kH_{E_{n+i}}(X,X)E_{n+i}+\nabla_XX\right)\\
  &=\sum_{i=1}^k\bigl(H_{E_{n+i}}(X,X)\widetilde\nabla_{\widetilde Y}E_{n+i}
    +\widetilde Y\cdot(H_{E_{n+i}}(X,X))E_{n+i}\bigr)
    +\widetilde\nabla_{\widetilde Y}\nabla_XX,
\end{align*}
有
\[
  \langle\widetilde\nabla_{\widetilde Y}\widetilde\nabla_{\widetilde X}\widetilde X,Y\rangle
  =\sum_{i=1}^k H_{E_{n+i}}(X,X)\langle\widetilde\nabla_{\widetilde Y}E_{n+i},Y\rangle
    +\langle\widetilde\nabla_{\widetilde Y}\nabla_XX,Y\rangle.
\]
此外,
\begin{align*}
  0&=\widetilde Y\cdot\langle E_{n+i},Y\rangle
    =\langle\widetilde\nabla_{\widetilde Y}E_{n+i},Y\rangle
    +\langle E_{n+i},\widetilde\nabla_{\widetilde Y}\widetilde Y\rangle\\
  &=\langle\widetilde\nabla_{\widetilde Y}E_{n+i},Y\rangle+\langle E_{n+i},B(Y,Y)+\nabla_YY\rangle\\
  &=\langle\widetilde\nabla_{\widetilde Y}E_{n+i},Y\rangle+\langle E_{n+i},B(Y,Y)\rangle\\
  &=\langle\widetilde\nabla_{\widetilde Y}E_{n+i},Y\rangle+H_{E_{n+i}}(Y,Y),
\end{align*}
所以
\begin{align*}
 &\langle\widetilde\nabla_{\widetilde Y}\widetilde\nabla_{\widetilde X}\widetilde X,Y\rangle\\
 &\qquad=-\sum_{i=1}^k H_{E_{n+i}}(X,X)H_{E_{n+i}}(Y,Y)
    +\langle\widetilde\nabla_{\widetilde Y}\nabla_XX,Y\rangle\\
 &\qquad=-\sum_{i=1}^k H_{E_{n+i}}(X,X)H_{E_{n+i}}(Y,Y)+\langle\nabla_Y\nabla_XX,Y\rangle.
\end{align*}
类似地, 可得
\[
  \langle\widetilde\nabla_{\widetilde X}\widetilde\nabla_{\widetilde Y}\widetilde X,Y\rangle
  =-\sum_{i=1}^k H_{E_{n+i}}(X,Y)H_{E_{n+i}}(X,Y)+\langle\nabla_X\nabla_YX,Y\rangle,
\]
从而
\begin{align*}
  K^N(\Pi)-K^M(\Pi)
  &=\sum_{i=1}^k\bigl(-(H_{E_{n+i}}(X_p,Y_p))^2+
     H_{E_{n+i}}(X_p,X_p)H_{E_{n+i}}(Y_p,Y_p)\bigr)\\
  &=-\lVert B(X_p,Y_p)\rVert^2+\langle B(X_p,X_p),B(Y_p,Y_p)\rangle.
\end{align*}
\end{RiemannProof}

\begin{RiemannExample}\label{ex:sectional-curvature-and-principal-curvatures}
仍考虑超曲面 $N$. 选取 $T_pN$ 的一个标准正交基 $\{(E_1)_p,\ldots,(E_n)_p\}$, 由 $S_{n_p}$ 的特征向量组成, 其中 $n_p\in(T_pN)^\perp$. 因此, 考虑 $T_pN$ 中由其中两个向量 $(E_i)_p,(E_j)_p$ 生成的截面 $\Pi$, 并利用 $B(X_p,Y_p)=\langle\!\langle S_{n_p}(X_p),Y_p\rangle\!\rangle n_p$, 有
\begin{align*}
  &K^N(\Pi)-K^M(\Pi)\\
  &\quad=-\lVert B((E_i)_p,(E_j)_p)\rVert^2
    +\langle B((E_i)_p,(E_i)_p),B((E_j)_p,(E_j)_p)\rangle\\
  &\quad=-\langle\!\langle S_{n_p}((E_i)_p),(E_j)_p\rangle\!\rangle^2
     +\langle\!\langle S_{n_p}((E_i)_p),(E_i)_p\rangle\!\rangle
      \langle\!\langle S_{n_p}((E_j)_p),(E_j)_p\rangle\!\rangle\\
  &\quad=\lambda_i\lambda_j.
\end{align*}
\end{RiemannExample}

\begin{RiemannExample}\label{ex:theorema-egregium}
在 $N$ 是 2 维流形且 $M=\mathbb{R}^3$ 带 Euclidean 度量的特殊情形下, 有 $K^M\equiv0$, 从而 $K^N(p)=\lambda_1\lambda_2$, 正如 Example~\ref{ex:principal-curvatures-of-hypersurface} 所述. 因此, 虽然 $\lambda_1$ 和 $\lambda_2$ 依赖于浸入, 但它们的\textbf{乘积}只依赖于 $N$ 的内在几何. Gauss 对这一发现十分满意, 将其称为他的 \textbf{Theorema Egregium} (\textquotedblleft 卓越定理\textquotedblright).
\end{RiemannExample}

现在详细研究 $N$ 是带 Euclidean 度量的 $M=\mathbb{R}^{n+1}$ 中的超曲面的特殊情形. 设 $c:I\to N$ 是 $N$ 中一条由弧长 $s$ 参数化的曲线, 满足 $c(0)=p$ 且 $\dot c(0)=X_p\in T_pN$. 将曲线 $c$ 与 $\mathbb{R}^{n+1}$ 中的曲线 $f\circ c$ 等同. 考虑定义在 $p$ 在 $N$ 中的邻域 $V$ 上的 Gauss 映射 $g:V\to S^n$, 在 $S^n$ 中取曲线 $n(s):=(g\circ c)(s)$. 由于 $\widetilde\nabla$ 是 $\mathbb{R}^3$ 上 Euclidean 度量对应的 Levi–Civita 联络, 有 $\langle\widetilde\nabla_{\dot c}\dot c,n\rangle=\langle\ddot c,n\rangle$. 另一方面,
\[
  \langle\widetilde\nabla_{\dot c}\dot c,n\rangle
  =\langle B(\dot c,\dot c)+\nabla_{\dot c}\dot c,n\rangle
  =\langle B(\dot c,\dot c),n\rangle=H_n(\dot c,\dot c).
\]
因此, 在 $s=0$ 时, $H_{g(p)}(X_p,X_p)=\langle\ddot c(0),n_p\rangle$. 这个值 $k_{n_p}:=\langle\ddot c(0),n_p\rangle$ 称为 $c$ 在 $p$ 处的\textbf{法曲率}. 由于 $k_{n_p}$ 等于 $H_{g(p)}(X_p,X_p)$, 它不依赖于曲线, 只依赖于曲线的初速度. 因为 $H_{g(p)}(X_p,X_p)=\langle\!\langle S_{g(p)}(X_p),X_p\rangle\!\rangle$, 在 $\lVert X_p\rVert=1$ 的约束下, 这些曲率的临界值等于 $\lambda_1,\ldots,\lambda_n$, 称为\textbf{主曲率}. 这就是在 Example~\ref{ex:principal-curvatures-of-hypersurface} 中也称 $S_{n_p}$ 的特征值为主曲率的原因. 于是 $f$ 的 Gauss 曲率等于主曲率的乘积, $K=\lambda_1\ldots\lambda_n$. 由于法曲率不依赖于在 $p$ 处与 $X_p$ 相切的曲线的选取, 可以取 $c$ 在 $X_p$ 与 $n_p$ 生成的 2 维平面内取值. 此时 $\ddot c(0)$ 平行于法向量 $n_p$, 所以
\[
  |k_n|=|\langle\ddot c(0),n\rangle|=\lVert\ddot c(0)\rVert=k_c,
\]
其中 $k_c:=\lVert\ddot c(0)\rVert$ 是曲线 $c$ 在 $c(0)$ 处的\textbf{曲率}. 若 $c$ 是 $N$ 的测地线, 同一公式也成立 [参见 Exercise~\ref{ex:isometric-immersions}(\ref{ex:principal-curvatures-from-geodesic-curvatures})].

\begin{RiemannExample}\label{ex:gauss-curvature-sphere-cylinder-saddle}
考虑下列三个曲面: 2 维球面, 圆柱面和鞍形曲面 $z=xy$.
\begin{enumerate}[label=(\arabic*),leftmargin=*]
\item 设 $p$ 是球面上任一点. 直观地, 这个曲面的所有点都在 $p$ 处切平面的同一侧, 这意味着两个主曲率同号 (取决于选取的定向), 因而所有点的 Gauss 曲率为正.
\item 若 $p$ 是圆柱面上任一点, 则其中一个主曲率为零 (是最大值还是最小值取决于所选定向), 所以所有点的 Gauss 曲率为零.
\item 最后, 若 $p$ 是鞍形曲面 $z=xy$ 上的一点, 则 $p$ 处的主曲率异号, 所以 Gauss 曲率为负.
\end{enumerate}
\end{RiemannExample}

\begin{RiemannExercise}\label{ex:isometric-immersions}
\begin{enumerate}[label=(\arabic*),ref=\arabic*,leftmargin=*]
\item\label{ex:second-fundamental-form-properties} 设 $M$ 是具有 Levi–Civita 联络 $\widetilde\nabla$ 的 Riemann 流形, $N$ 是赋有诱导度量和 Levi–Civita 联络 $\nabla$ 的子流形. 设 $\widetilde X,\widetilde Y\in\mathfrak{X}(M)$ 为 $X,Y\in\mathfrak{X}(N)$ 的局部延拓. 回顾, $N$ 在 $M$ 中的包含映射的第二基本形式是每点 $p\in N$ 处由下式定义的映射 $B:T_pN\times T_pN\to(T_pN)^\perp$:
\[
  B(X,Y):=\widetilde\nabla_{\widetilde X}\widetilde Y-\nabla_XY.
\]
证明:
\begin{enumerate}[label=(\alph*),leftmargin=*]
\item $B(X,Y)$ 不依赖于延拓 $\widetilde X,\widetilde Y$ 的选取;
\item $B(X,Y)$ 垂直于 $N$;
\item $B$ 对称, 即 $B(X,Y)=B(Y,X)$;
\item $B$ 双线性;
\item $B(X,Y)_p$ 只依赖于 $X_p$ 和 $Y_p$ 的取值;
\item $\widetilde\nabla_{[\widetilde X,\widetilde Y]}\widetilde X-\nabla_{[X,Y]}X$ 垂直于 $N$.
\end{enumerate}
\item\label{ex:constant-curvature-of-round-sphere} 设 $S^n(r)\subset\mathbb{R}^{n+1}$ 是半径为 $r$ 的 $n$ 维球面.
\begin{enumerate}[label=(\alph*),leftmargin=*]
\item 在每点选取向外的单位法向量, 这个包含映射的 Gauss 映射是什么?
\item 它的导数的特征值是什么?
\item 证明所有截面曲率都等于 $\frac1{r^2}$ (所以 $S^n(r)$ 的曲率恒为 $\frac1{r^2}$).
\end{enumerate}
\item\label{ex:totally-geodesic-submanifolds} 设 $(M,\langle\cdot,\cdot\rangle)$ 是 Riemann 流形. 若 $N$ 的测地线是 $M$ 的测地线, 则称子流形 $N\subset M$ 是\textbf{全测地的}. 证明:
\begin{enumerate}[label=(\alph*),leftmargin=*]
\item $N$ 全测地当且仅当 $B\equiv0$, 其中 $B$ 是 $N$ 的第二基本形式;
\item 若 $N$ 是某个等距映射的不动点集, 则 $N$ 全测地. 利用该结果给出 $\mathbb{R}^n$, $S^n$ 和 $H^n$ 的全测地子流形的例子.
\end{enumerate}
\item\label{ex:gauss-curvature-as-ratio-of-volumes} 设 $N$ 是 $\mathbb{R}^{n+1}$ 中的超曲面, $p$ 是 $N$ 中一点. 证明若 $K(p)\ne0$, 则
\[
  |K(p)|=\lim_{D\to p}\frac{\operatorname{vol}(g(D))}{\operatorname{vol}(D)},
\]
其中 $g:V\subset N\to S^n$ 是 Gauss 映射, $D$ 是 $p$ 的邻域, 其直径趋于零.
\item\label{ex:sectional-curvature-by-geodesic-surface} 设 $(M,\langle\cdot,\cdot\rangle)$ 是 Riemann 流形, $p$ 是 $M$ 中一点, $\Pi$ 是 $T_pM$ 的一个截面. 对 $p$ 附近的正规球 $B_\varepsilon(p):=\exp_p(B_\varepsilon(0))$, 考虑集合 $N_p:=\exp_p(B_\varepsilon(0)\cap\Pi)$. 证明:
\begin{enumerate}[label=(\alph*),leftmargin=*]
\item 集合 $N_p$ 是 $M$ 的 2 维子流形, 由 $B_\varepsilon(p)$ 内在 $p$ 处与 $\Pi$ 相切的测地线段构成;
\item 若在 $N_p$ 上使用由 $M$ 的度量诱导的度量, 则截面曲率 $K^M(\Pi)$ 等于 2 维流形 $N_p$ 的 Gauss 曲率.
\end{enumerate}
\item\label{ex:principal-curvatures-from-geodesic-curvatures} 设 $(M,\langle\cdot,\cdot\rangle)$ 是具有 Levi–Civita 联络 $\widetilde\nabla$ 的 Riemann 流形, $N$ 是 $M$ 中的一个超曲面. 由弧长参数化的曲线 $c:I\subset\mathbb{R}\to M$ 的\textbf{测地曲率}为 $k_g(s)=\lVert\widetilde\nabla_{\dot c(s)}\dot c(s)\rVert$. 证明主曲率的绝对值是与主方向相切的 $N$ 的测地线 (在 $M$ 中) 的测地曲率.
{\small (\textbf{Remark:} 在定向 2 维 Riemann 流形的情形, 根据 $\{\dot c(s),\widetilde\nabla_{\dot c(s)}\dot c(s)\}$ 的定向, $k_g$ 取正或负, 参见第~\ref{sec:cartan-structure-equations}~节).}
\item\label{ex:gauss-maps-of-paraboloid-and-saddle} 利用 Gauss 映射计算 $\mathbb{R}^3$ 中下列曲面的 Gauss 曲率:
\begin{enumerate}[label=(\alph*),leftmargin=*]
\item 抛物面 $z=\frac12(x^2+y^2)$;
\item 鞍形曲面 $z=xy$.
\end{enumerate}
\item\label{ex:fundamental-forms-of-surfaces-of-revolution} (\emph{旋转曲面}) 考虑由
\[
  f(s,\theta)=(h(s)\cos\theta,h(s)\sin\theta,g(s))
\]
给出的映射 $f:\mathbb{R}\times(0,2\pi)\to\mathbb{R}^3$, 其中 $h>0$ 和 $g$ 为满足
\[
  (h'(s))^2+(g'(s))^2=1
\]
的光滑映射. $f$ 的像是以 $Oz$ 为轴的旋转曲面 $S$, 由弧长 $s$ 参数化的曲线 $\alpha(s)=(h(s),g(s))$ 绕此轴旋转而成.
\begin{enumerate}[label=(\alph*),leftmargin=*]
\item 证明 $f$ 是浸入.
\item 证明 $f_s:=(\mathrm{d}f)\bigl(\frac{\partial}{\partial s}\bigr)$ 和 $f_\theta:=(\mathrm{d}f)\bigl(\frac{\partial}{\partial\theta}\bigr)$ 正交.
\item 求 Gauss 映射并计算 $S$ 相对于标架 $\{E_s,E_\theta\}$ 的第二基本形式的矩阵, 其中 $E_s:=f_s$, $E_\theta:=\frac1{\lVert f_\theta\rVert}f_\theta$.
\item 计算 $S$ 的平均曲率 $H$ 和 Gauss 曲率 $K$.
\item 利用这些结果, 分别给出满足以下条件的旋转曲面的例子:
\begin{enumerate}[label=(\arabic*),leftmargin=*]
\item $K\equiv0$;
\item $K\equiv1$;
\item $K\equiv-1$;
\item $H\equiv0$ (不是平面).
\end{enumerate}
{\small (\textbf{Remark:} 平均曲率恒为零的曲面称为\textbf{极小曲面}. 可以证明, 若一个紧带边界曲面在所有具有相同边界的曲面中面积最小, 则它必为极小曲面).}
\end{enumerate}
\end{enumerate}
\end{RiemannExercise}
